\documentclass[11pt]{article}
\usepackage[T1]{fontenc}
\usepackage{lmodern,amsmath,amssymb,amsthm,mathtools,microtype}
\usepackage[a4paper,margin=27mm,top=24mm,bottom=24mm]{geometry}
\usepackage{xcolor,hyperref}
\definecolor{ink}{HTML}{163C58}
\hypersetup{colorlinks=true,linkcolor=ink,urlcolor=ink,citecolor=ink,pdftitle={Linear relations and algebraic independence for weighted Fibonacci and Lucas sums}}
\newtheorem{theorem}{Theorem}
\newtheorem{lemma}[theorem]{Lemma}
\newcommand{\QQ}{\mathbb Q}
\newcommand{\CC}{\mathbb C}
\newcommand{\F}{\mathcal F}
\newcommand{\coef}[1]{[q^{#1}]}
\setlength{\parindent}{1em}
\setlength{\emergencystretch}{1.4em}
\pagestyle{plain}
\begin{document}
\begin{center}
{\LARGE\bfseries Linear relations and algebraic independence\\[2pt] for weighted Fibonacci and Lucas sums}\par
\vspace{0.6em}
{\large A polynomial-weight extension of Stein's conjecture}\par
\vspace{0.7em}
Revised research note \quad $\cdot$ \quad 15 November 2026
\end{center}
\smallskip
\noindent\textbf{Abstract.}
For polynomially weighted even reciprocal powers of Fibonacci and Lucas numbers, including alternating sums, we determine the entire space of linear relations over the algebraic numbers. The rank depends on the interaction of the weight space with the dilation $P(X)\mapsto P(4X)$. A single monomial yields rank $3M$ through denominator exponent $2M$. A polynomial with two or more terms yields rank $4M$, and every nontrivial algebraic linear combination is transcendental. This extends the unweighted theorem of Elsner and Technau on Stein's conjecture. For any three distinct positive indices $s_i$ and any polynomials $W_i$ of the same positive degree with algebraic coefficients, we prove that the three ordinary Fibonacci sums with terms $W_i(n^2)/F_n^{2s_i}$ are algebraically independent. The same assertion holds for algebraic-root Lucas sequences of the first kind with root product $-1$ and one root strictly inside the unit disk.

\section{The arithmetic statement}
Let $F_1=F_2=1$ and $F_{n+2}=F_{n+1}+F_n$. All series below converge absolutely by exponential growth.
\begin{theorem}\label{main}
For every $1\le s_1<s_2<s_3$, the numbers
\[
 S_{s_i}=\sum_{n=1}^{\infty}\frac{n^2}{F_n^{2s_i}}
 \qquad (i=1,2,3)
\]
are algebraically independent over $\QQ$. Any four distinct members of the sequence satisfy a nonzero polynomial relation over $\QQ$. Thus any three distinct members form a transcendence basis for $\QQ(S_s:s\ge1)$.
\end{theorem}
Wang's published article \cite{Wang2025} treats several families of hyperbolic and recurrence sums. The point of comparison here is specifically Proposition~5.10 in its second arXiv version. Its argument displays the relevant highest-weight forms and imposes a finite upper bound in testing their independence. We use the same reduction to quasimodular forms and prove independence with different Fourier coefficients. The polynomial-weight rank formula appears in Theorem~\ref{rank}, the quartic numerator result in Theorem~\ref{quarticthm}, and the unrestricted common-degree theorem in Theorem~\ref{full}.

\subsection*{The specialization used below}
Write $q=e^{2\pi iz}$ and $D=q\,d/dq$. Wang's specialization theorem \cite[Theorem 3.5]{Wang2024}, also stated as Theorem~5.3 in \cite{Wang2025}, applies to finite algebraic linear combinations of quasimodular forms with algebraic Fourier coefficients. If three such functions are algebraically independent over $\CC$, their values at an algebraic $q$ with $0<|q|<1$ are algebraically independent over $\QQ$. For a finite sum of homogeneous forms we call its nonzero component of largest weight its top component. If the three top components are algebraically independent, so are the original functions. Indeed, a polynomial relation would give a nonzero relation among the top components by taking the part of maximal weighted degree \cite[Lemma 5.4]{Wang2025}.

\clearpage
\section{From Fibonacci numbers to the top components}
Set $\phi=(1+\sqrt5)/2$ and $q_0=\phi^{-2}$. Binet's formula and the binomial series give
\begin{align}
 S_s&=5^s\F_s(q_0),\nonumber\\
 \F_s(q)&=\sum_{n,m\ge1}n^2 p_s(m)(-1)^{n(m-s)}q^{nm},&
 p_s(m)&=\frac{m\prod_{j=1}^{s-1}(m^2-j^2)}{(2s-1)!}.\label{raw}
\end{align}
In fact $p_s(m)=\binom{m+s-1}{2s-1}$ for $m\ge s$ and vanishes for $1\le m<s$. All rearrangements in \eqref{raw} are absolutely convergent for $|q|<1$.

Use the normalized Eisenstein series $E_{2k}(z)=1+a_k\sum_{m,n\ge1}m^{2k-1}q^{mn}$, where $a_k=-4k/B_{2k}\ne0$ and $B_j$ denotes the $j$th Bernoulli number. This convention includes $E_2$. Define
\begin{align}
 g(z)&=E_2(z)-2E_2(2z)+4E_2(4z),\nonumber\\
 f_k(z)&=E_{2k}(z)-\bigl(4^k+1-(-1)^k\bigr)E_{2k}(2z)
          +4^k E_{2k}(4z) \qquad(k\ge1).\label{fk}
\end{align}
Up to nonzero rational factors, the top components of $\F_s$ are
\begin{equation}\label{top}
 T_1=Dg,\qquad T_{k+1}=D^2 f_k\quad(k\ge1).
\end{equation}
For completeness, in \eqref{raw} a monomial $m^{2j-1}$ with $j\ge2$ gives $D^2$ applied to the signed Lambert sum with $m^{2j-3}$, producing weight $2j+2$. For the top term $j=s$, if $s$ is even the sign is $(-1)^{nm}$. If $s$ is odd the sign is $-(-1)^{nm}+2\boldsymbol{1}_{2\mid n}$. Splitting the two parity classes gives \eqref{fk}, up to a scalar and an additive constant removed by $D^2$. Equivalently, one uses the divisor-sum identity
\[
 E_{2k}(z+\tfrac12)=-E_{2k}(z)+(2+4^k)E_{2k}(2z)-4^kE_{2k}(4z).
\]
For $j=1$, take one $D$ and interchange the indices to obtain weight four. When $s=1$, the signed sum after this interchange is $-E_2(z+1/2)+4E_2(2z)$ up to scaling and a constant. This equals $g(z)$ up to that normalization. Lower powers in $p_s$ have strictly smaller weights. These formulas agree with the top components displayed in \cite[Proposition 5.10]{Wang2025}.

\clearpage
\section{A determinant for functional independence}
Use the derivations $D,L,W$ on the graded quasimodular ring of level $4$. Here $W$ multiplies a homogeneous form by its weight, $L$ kills modular forms, $LE_2=12$, and $[L,D]=W$, $[W,D]=2D$ \cite[Section 8]{NSZ}. In particular, whenever $f$ is homogeneous of weight $w$,
\[
 LDf=D Lf+wf,\qquad LD^2f=D^2 Lf+2(w+1)Df.
\]
The forms $f_k$ are modular for $k\ge2$. For $k=1$ and for $g$, the identity $mE_2(mz)-E_2(z)\in M_2(\Gamma_0(m))$ gives $Lf_1=-12$ and $Lg=12$. Constancy of these two derivatives is enough.

For three functions $T_i$ form the three-row matrix with $i$th row
\[
 (D T_i,\, L T_i,\, W T_i).
\]
A nonzero determinant implies algebraic independence over $\CC$. Indeed, differentiating a hypothetical relation of minimal total degree gives the gradient as a left null vector of the matrix. Some partial derivative of the relation is a nonzero polynomial. Its value at the three functions is nonzero by minimality. Hence a relation forces determinant zero.

\subsection*{Case I. All indices are at least two}
For $T_{k+1}=D^2 f_k$ the row is exactly
\[
 (D^3 f_k,\,2(2k+1)Df_k,\,(2k+4)D^2 f_k).
\]
In particular the same row formula works for $k=1$. From \eqref{fk}, $[q]f_k=a_k$ and
\[
 \rho_k=\frac{[q^2]f_k}{[q]f_k}=(-1)^k-\frac{4^k}{2}.
\]
The row divided by $a_k$ has expansion $qv(k)+q^2\rho_k u(k)+O(q^3)$, with
\[
 v(k)=(1,4k+2,2k+4),\qquad u(k)=(8,8k+4,8k+16).
\]
Write $v(k)=v_0+kv_1$, and let $n=(6,1,-2)$. Then $v_0\times v_1=-2n$. All $v(k)$ lie in the plane perpendicular to $n$, so the coefficient of $q^3$ in the determinant vanishes. By multilinearity the coefficient of $q^4$ in the determinant of the normalized rows is
\begin{equation}\label{cert}
 -2\det\begin{pmatrix}
 1&a&h(a)\\1&b&h(b)\\1&c&h(c)
 \end{pmatrix},\qquad h(k)=\rho_k\, n\!\cdot\!u(k)=(20-8k)\left((-1)^k-\frac{4^k}{2}\right)
\end{equation}
for indices $a+1<b+1<c+1$. For example, decompose each $\rho_k u(k)$ into its component in the plane of $v_0,v_1$ and a multiple of a fixed vector $e$ with $n\cdot e=1$. Only the latter component contributes, with coefficient $h(k)$, and $\det(v_0,v_1,e)=-2$.

The forward second difference is
\[
 h(k+2)-2h(k+1)+h(k)
 =4^k(36k+6)-(-1)^k(32k-48)>0\qquad(k\ge1).
\]
For $k=1$ the value is $152$. For $k\ge2$ the first term exceeds $|32k-48|$. Thus the consecutive slopes increase strictly. The secant slope on $[a,b]$ is smaller than the secant slope on $[b,c]$ for all $a<b<c$. This makes the determinant in \eqref{cert} nonzero for every triple.

\clearpage
\subsection*{Case II. The smallest index is one}
For $T_1=Dg$, the row is $(D^2g,2g,4Dg)$. Its constant term is $(0,6,0)$ because $g$ has constant term $3$. For any $1\le a<b$, the determinant with the next two rows normalized by $a_a$ and $a_b$ has $q^2$ coefficient
\[
 \det\begin{pmatrix}
 0&6&0\\1&4a+2&2a+4\\1&4b+2&2b+4
 \end{pmatrix}=-12(b-a)\ne0.
\]
This covers $(1,2,s)$ as well. Both cases in \eqref{top} are settled. The top-weight argument and Wang's specialization theorem at the algebraic $q_0$ now prove the first assertion of Theorem~\ref{main}.

\section{The upper bound and a check on the exceptional triple}
At fixed level, the dimension of the space of modular forms of weight $j$ is $O(j+1)$ by the valence bound. Quasimodular forms are polynomials in $E_2$ with modular coefficients. Consequently the space of sums of quasimodular forms of weights up to $K$ has dimension $O(K^3)$. Four fixed functions $\F_s$ have $\binom{d+4}{4}$ monomials of total degree at most $d$, all in weights $O(d)$. For large $d$ this forces a complex linear relation among the monomials. Their rational Fourier coefficients imply a rational relation: the infinite coefficient matrix has finitely many columns, and its rational row rank equals its complex row rank, as is seen by its finite minors. Evaluation at $q_0$ and the rational factors $5^s$ give the required dependence of four numbers. This completes Theorem~\ref{main}.

\paragraph{The triple $(4,6,8)$.}
In the notation of \eqref{fk}, $f_k=-E_{2k}(z+1/2)$ for every odd $k$. The classical identities $E_{10}=E_4E_6$ and $E_{14}=E_4^2E_6$ therefore yield
\[
 f_3 f_7=f_5^2.
\]
Thus the quotient $f_3^2f_7^2/f_5^4$ equals one identically. The nonzero $q^2$ coefficient stated for this quotient in the public v2 proof of Proposition~5.10 is incompatible with the displayed forms (our subscript $k$ is Wang's $s-1$). The differentiated forms have a different outcome. In \eqref{cert}, $(a,b,c)=(3,5,7)$ gives the nonzero normalized coefficient $-1{,}098{,}240$. This proves the numerical independence assertion for the exceptional triple within the uniform argument.

\clearpage
\section{The polynomial-weight version of Stein's conjecture}\label{stein}
Put $L_0=2,L_1=1,L_{n+2}=L_{n+1}+L_n$, and $K=\overline{\QQ}$. For $Z=F,L$, $\varepsilon=+1,-1$, and a polynomial $P\in K[X]$, define
\[
 Z_s^\varepsilon[P]=\sum_{n\ge1}\frac{\varepsilon^n P(n^2)}{Z_n^{2s}}.
\]
Our minus sign uses $(-1)^n$. The starred convention of Stein and of Elsner--Technau is its negative. Define the invertible linear map $T$ on $K[X]$ by $(TP)(X)=P(4X)$.

\begin{theorem}[Exact rank for polynomial weights]\label{rank}
Let $M\ge1$, let $V\subseteq K[X]$ be finite dimensional with basis $P_1,\ldots,P_b$, and put $d_V=\dim_K\{B\in V:TB\in V\}$. The $4Mb$ numbers $Z_s^\varepsilon[P_i]$, $1\le s\le M$, have
\[
 \boxed{\quad\text{linear rank }4Mb-Md_V,\qquad
                 \text{relation dimension }Md_V.\quad}
\]
Their span intersects $K$ only in zero. For a rational weight space the same rank formula holds over $\QQ$.
\end{theorem}
\begin{theorem}[One polynomial]\label{dich}
If $P\ne0$, the rank for the four families through $s=M$ is $3M$ when $P$ is a monomial and $4M$ when $P$ has at least two nonzero coefficients. In the second case every nontrivial finite $K$-linear combination of these four infinite families is transcendental.
\end{theorem}
For $P=1$, Stein conjectured that $M$ independent relations exhaust the kernel, together with a parity restriction on the Lucas coefficients. Elsner and Technau proved both assertions \cite[Conjecture and Theorem 1.1]{ET}. Theorem~\ref{rank} gives a weighted extension and classifies relations between different numerator weights as well. The partial fractions below also recover the parity restriction. For example, $P=1+X$ eliminates every nonzero relation for every $M$. The vector space $V=\langle 1,X\rangle$, by contrast, has rank $6M$ on $8M$ values. A space with an intermediate answer is $V=\langle1, X+X^2\rangle$. Here $d_V=1$, so the rank is $7M$.

\subsection*{Specializing linear relations}
A nonconstant mixed-weight quasimodular form $G$ of level four with algebraic Fourier coefficients takes a transcendental value at every algebraic $q$ with $0<|q|<1$. Indeed, $G$ is transcendental over the constant field $\CC$. The independent functions $E_2,E_4,E_6$ give a transcendence basis for the level-four quasimodular field: independence holds at level one \cite{NSZ}, and the dimension argument of the preceding section gives an upper bound of three. By basis exchange, $G$ can be supplemented by two of these functions to form an algebraically independent triple. Wang's specialization theorem then applies. In particular, on a span of such forms with zero constant Fourier coefficient, a combination with algebraic coefficients takes an algebraic value at $q_0$ exactly when it vanishes identically.

\clearpage
\section{Independent Lambert series and all relations}\label{linearproof}
For $u\ge0$ even, $v\ge1$ odd and $e\in\{1,2,4\}$ write
\[
 H_{u,v}(q^e)=\sum_{n,m\ge1}n^u m^v q^{enm}.
\]
These functions have zero constant term and are mixed-weight quasimodular forms with rational Fourier coefficients. If $a=\min(u,v)$, $k=|u-v|+1$, then $k\ge2$ is even. The identity $n^um^v=(nm)^a n^{u-a}m^{v-a}$ and symmetry in $n,m$ express $H_{u,v}(q^e)$ as a nonzero rational multiple of $D^a(E_k(ez)-1)$. The factor from the substitution $q\mapsto q^e$ is $e^{-a}$.

\begin{lemma}\label{prime}
The series $H_{u,v}(q^e)$ indexed by these triples $(u,v,e)$ are linearly independent over $\CC$. Their values at $q_0$ together with $1$ are linearly independent over $K$.
\end{lemma}
\begin{proof}
In a putative finite linear relation examine the coefficient at $q^{\ell r}$ for distinct odd primes $\ell,r$. Only $e=1$ contributes. The coefficient of $H_{u,v}(q)$ is
\[
 (\ell^u+\ell^v)(r^u+r^v).
\]
A bivariate polynomial vanishing at all such pairs vanishes identically: hold one prime fixed and vary the other, then vary the first. In the resulting zero polynomial the coefficient of $X^uY^v$ is precisely the coefficient of $H_{u,v}(q)$ in the relation, since $u$ is even and $v$ is odd. With all $e=1$ terms removed, apply the same comparison at $q^{2\ell r}$, then at $q^{4\ell r}$. This proves functional independence. The preceding specialization argument and the zero constant terms prove the second assertion.
\end{proof}

\paragraph{A rectangular family.}
Set $A_s^\varepsilon[P]=5^{-s}F_s^\varepsilon[P]$, $B_s^\varepsilon[P]=L_s^\varepsilon[P]$. For $P=X^p$, Binet's formula expresses these four functions before specialization as the double sum in \eqref{raw} with $n^{2p}$ in place of $n^2$ and with the four multipliers
\[
 (-1)^{n(m-s)},\quad (-1)^n(-1)^{n(m-s)},\quad
 (-1)^{(n+1)(m-s)},\quad (-1)^n(-1)^{(n+1)(m-s)} .
\]
These span all four functions of $n,m$ modulo two. Adding and subtracting the twists separates even and odd $n$. On either parity of $n$, one denominator then gives a constant in $m$ and the other gives a nonzero multiple of $(-1)^m$.

The polynomials $p_s(m)$ of \eqref{raw}, $1\le s\le M$, form a basis of the odd polynomials of degree at most $2M-1$, by their distinct degrees. Combining this triangular change of basis with the parity filters gives precisely the span of
\[
 H_{2p,2j-1}(q^e),\qquad 1\le j\le M,\quad e=1,2,4 .
\]
Indeed, the unrestricted sum gives $e=1$, restriction to even $n$ or even $m$ gives a scalar multiple of $e=2$, and restriction to both even gives $e=4$. Consequently for $p=0,\ldots,P_0$ the full rectangular array has rank $3M(P_0+1)$ by Lemma~\ref{prime}. Its relation spaces at the distinct $p$ form a direct sum.

\clearpage
\subsection*{An explicit basis of identities}
Put $x_n=5F_n^2$. The duplication identities are
\[
 L_n^2=x_n+4(-1)^n,\qquad x_{2n}=x_n L_n^2 .
\]
For $1\le j\le s$ set $c_{sj}=(-1)^{s-j}\binom{2s-j-1}{s-1}4^{j-2s}$ and write
\[
 C_s[P]=\sum_{j=1}^s c_{sj}
 \bigl(A_j^{(-1)^j}[P]+(-1)^jB_j^{(-1)^j}[P]\bigr).
\]
Then for every polynomial $P$ the following identity holds:
\begin{equation}\label{allrelations}
 \frac{A_s^+[P]+A_s^-[P]}{2}= C_s[TP].
\end{equation}
To see it, the left side is the sum over even indices, whose summand at $2n$ is $P(4n^2)/x_{2n}^s$. Expand
\[
 \frac{1}{x^s(x+a)^s}=\sum_{j=1}^s
 (-1)^{s-j}\binom{2s-j-1}{s-1}a^{j-2s}
 \left(\frac1{x^j}+\frac{(-1)^j}{(x+a)^j}\right).
\]
The principal parts at $x=0,-a$ give this identity and the difference tends to zero at infinity. Put $x=x_n$, $a=4(-1)^n$ and sum.

For $P=X^p$ we have $TP=4^pP$. Formula~\eqref{allrelations} supplies $M$ linearly independent relations: in $C_s$ the Lucas coordinate with exponent $j=s$ has nonzero coefficient, and larger exponents are absent. The rectangular rank above proves these exhaust the kernel. The displayed basis uses only the alternating Lucas coordinate for odd $j$ and only the ordinary Lucas coordinate for even $j$.

\paragraph{Restricting the weights.}
For an arbitrary finite collection of polynomial coefficients, choose $P_0$ bounding all degrees and expand them in monomials. The rectangular result shows that a relation in the normalized coordinates has a unique representation as the sum over $s$ of the left sides minus the right sides of \eqref{allrelations}, with individual parameter polynomials $P_s$. If all coordinate polynomials belong to $V$, the triangular Lucas coefficients, read starting at $s=M$, imply $TP_s\in V$ for every $s$. The Fibonacci coordinates then imply $P_s\in V$, since their remaining terms are just $P_s/2$. Conversely these two conditions suffice. There are $M$ independent choices of $P_s\in V\cap T^{-1}V$. The normalization by $5^{-s}$ is rational and invertible. This proves Theorem~\ref{rank}. All linear algebra for a rational weight space is defined over $\QQ$, so the rank is unchanged after extension to $K$. The absence of nonzero algebraic values in the span follows from Lemma~\ref{prime}.

When $V=KP$, the condition $TP\in KP$ is the eigenvector equation for $T$. The monomials $X^p$ have pairwise distinct eigenvalues $4^p$. This proves Theorem~\ref{dich}. For example at $s=1$, write $F^*,L^*$ using $(-1)^{n+1}$ and suppress the index $s$. For $P=X^p$, \eqref{allrelations} yields
\[
 F^+[P]+(2^{2p-1}-1)F^*[P]-5\,2^{2p-1}L^*[P]=0.
\]
At $p=0$ this recovers the motivating identity in \cite{ET}. At $P=1+X$, the four series together with $1$ are linearly independent over $K$.

\clearpage
\section{Algebraic independence for quartic numerator weights}\label{quartic}
\begin{theorem}\label{quarticthm}
Let $W_1,W_2,W_3\in K[X]$ all have degree two. For every $1\le s_1<s_2<s_3$ the numbers $F_{s_i}^+[W_i]$ are algebraically independent. The same conclusion holds if the three weights all have degree one.
\end{theorem}
The numerator of $F_s^+[W]$ is $W(n^2)$. Thus degree two includes $n^4$, with arbitrary lower even powers and a nonzero leading coefficient. This is algebraic independence for three ordinary Fibonacci sums. Theorems~\ref{rank} and \ref{dich} concern linear relations for all four types of sums.

To determine the top component, replace $n^2$ by $n^{2p}$ in \eqref{raw}. Each monomial $n^{2p}m^{2j-1}$ leads to derivatives of an Eisenstein series. Put
\[
 d=\min(2p,2j-1),\qquad w=|2p-(2j-1)|+1 .
\]
Indeed, strip the factor $(nm)^d$ via $D^d$, then impose the parity filters as in Section~\ref{linearproof}. The resulting weight is $w+2d=2(p+j)$, apart from a possible constant term for $d=0$. Thus only the leading monomial $n^{2p}m^{2s-1}$ contributes at weight $2(p+s)$. This proves the degree-one assertion from Theorem~\ref{main} and reduces degree two to $p=2$.

For $p=2$ write the top component, up to a nonzero rational factor, as $D^d G_s$. Use $E_w(tz)$ in the order $t=1,2,4$ in the last column:
\[
\begin{array}{c|c|c|c}
 s & d & w &\text{coefficients of }G_s\\ \hline
 1 &1&4&(1,-2,16)\\
 2 &3&2&(1,-6,4)\\
 s\ge3 &4&2(s-2)&(1,-(4^{s-2}+1+(-1)^s),4^{s-2})
\end{array}
\]
For the last row the factor remaining after $D^4$ is $m^{2s-5}$; separate even and odd indices as in \eqref{fk}. In the first row the remaining factor is $n^3$, and the sign $(-1)^{n(m-1)}$ splits as $-(-1)^{nm}+2\boldsymbol1_{2\mid n}$. This yields the stated $(1,-2,16)$. In the second row the remaining factor is $n$ with sign $(-1)^{nm}$, giving $-E_2(z+1/2)$ up to scaling. Each $LG_s$ is constant, including in weight two.

\paragraph{The differential determinant.}
Use the same three columns $D,L,W$ as in Section~3 (in these columns $W$ denotes the weight operator). The commutator identity
\[
 LD^d G_s=D^d LG_s+d(w+d-1)D^{d-1}G_s
\]
follows by induction. For $s\ge2$ normalize each row by dividing by $[q]G_s\ne0$, and write $\rho_s=[q^2]G_s/[q]G_s$. The first coefficient is the vector
\[
 v_s=(1,4(2s-1),2s+4)
\]
and the second coefficient is $\rho_s\, u_s$, where
\[
 u_s=\bigl(2^{d+1},\ 2^{d-1}4(2s-1),\ 2^d(2s+4)\bigr).
\]
Here $d(w+d-1)=4(2s-1)$ on both sides of $s=2$. A normal to the plane of the $v_s$ is $N=(20,1,-4)$, and $N\cdot u_s=4\cdot 2^{d-1}(11-2s)$. As in \eqref{cert}, the $q^4$ coefficient of the determinant is nonzero whenever the three points $(s,H_s)$ are not collinear, where
\[
 H_s=2^{d-1}(11-2s)\rho_s.
\]

\clearpage
\paragraph{The secant slopes.}
The table of $G_s$ gives $H_2=-84$. For $k=s-2\ge1$ it gives
\begin{equation}\label{quartich}
 H_{k+2}=8h_k,\qquad h_k=(2k-7)\bigl(4^k/2+(-1)^k\bigr).
\end{equation}
In particular
\[
 (H_2,H_3,H_4,H_5,H_6)=(-84,-40,-216,-248,1032).
\]
The forward second difference of $h_k$ equals
\[
 (9k-15/2)4^k+(8k-20)(-1)^k>0.
\]
At $k=1$ it is $18$, and at $k\ge2$ the exponential term dominates the absolute value of the second term. Consequently all triples with $s_i\ge3$ are noncollinear. For triples starting at $s=2$, the four secant slopes to $s=3,4,5,6$ are $44,-66,-164/3,279$. They are pairwise distinct. The consecutive slopes from $s=5$ onward are strictly increasing, and $H_6-H_5=1280>279$. Every subsequent secant slope from 2 is a weighted average of the previous slope from 2 and the next consecutive slope. It therefore increases strictly from $279$ onward. This covers every triple starting at 2.

For $s=1$, the row of $DG_1$ has constant term $(0,60,0)$ since $G_1(0)$ in the $q$-coordinate is $15$ and $w=4$. The other two rows, indexed by $a<b$, start at $qv_a$ and $qv_b$. The coefficient at $q^2$ is $-120(b-a)\ne0$. The determinant now proves independence in every case, and specialization finishes the proof of Theorem~\ref{quarticthm}.

\paragraph{An infinite independent family in a field of transcendence degree three.}
Consider the four signed families for $P=1+X$, with all $s\ge1$. By Theorem~\ref{dich} every nontrivial finite algebraic linear combination is transcendental. By Theorem~\ref{quarticthm} any three ordinary Fibonacci members are algebraically independent. All four families belong to the field over $K$ generated by the values of level-four quasimodular forms with algebraic Fourier coefficients, of transcendence degree at most three, by the dimension argument in Section~4. Thus adjoining the Lucas and alternating values does not increase the transcendence degree, even though each one supplies a further linearly independent element.

\section{Three modular derivatives of a common order}\label{high}
The following argument also gives an infinite range of algebraic independence for every higher numerator degree. A rising factorial is denoted by $(k)_j$ with $(k)_0=1$.

\begin{lemma}\label{deriv}
Suppose $f_i=a_i+c_iq+O(q^2)$ are modular forms of distinct positive integral weights $k_i$ on a common congruence subgroup, with $a_ic_i\ne0$ for $i=1,2,3$. For every $r\ge2$ the functions $D^rf_i$ are algebraically independent over $\CC$.
\end{lemma}
\begin{proof}
On the quasimodular ring use the homomorphism sending a homogeneous form $G$ of weight $w$ to $t^w\exp(XL)G$. The exponential is a finite sum on each form since $L$ lowers weight. The operator is multiplicative because $L$ is a derivation. The commutator formula above and $Lf_i=0$ give
\[
 \exp(XL)D^rf_i=
 \sum_{j=0}^r \binom rj(k_i+j)_{r-j}X^{r-j}D^j f_i.
\]
Indeed iterated lowering of $D^rf_i$ multiplies by $r(k_i+r-1)$ at its first step, then $(r-1)(k_i+r-2)$, and so on.

\clearpage
Apply the derivatives $t\partial_t$, $X\partial_X$, and $\partial_q$ to the three images and set $q=0$. Dividing each row by $a_i(k_i)_r t^{k_i+2r}X^r$ yields
\[
 \left(k_i+2r,\ r,\ \frac{c_i}{a_i}P_{k_i,r}(1/X)\right),
 \qquad
 P_{k,r}(Y)=\sum_{j=0}^r\binom rj\frac{Y^j}{(k)_j}.
\]
The three polynomials are linearly independent. After multiplication by $k_i(k_i+1)$, their coefficients in degrees $0,1,2$ are the values at $k_i$ of $k(k+1),\ r(k+1),\ \binom r2$; their evaluation matrix is invertible at three distinct $k_i$. The cofactors of the last column in the displayed Jacobian are nonzero since the weights are distinct. Its determinant is therefore a nonzero polynomial in $1/X$. At a generic nonzero $X,t$ the analytic inverse function theorem gives an open set of values of the three images, precluding any polynomial relation. A relation among the original functions would persist under the homomorphism, so it too is impossible.
\end{proof}

\begin{theorem}\label{upperthm}
Let $p\ge1$ and let $W_1,W_2,W_3\in K[X]$ all have degree $p$. For any $p+2\le s_1<s_2<s_3$, the three numbers $F_{s_i}^+[W_i]$ are algebraically independent. For $p=1,2$, Theorem~\ref{quarticthm} gives the full positive range.
\end{theorem}
\begin{proof}
For each $s\ge p+2$ set $w=2(s-p)\ge4$. The top component is proportional to $D^{2p}G_s$, now with
\[
 G_s=E_w(z)-(2^w+1+(-1)^s)E_w(2z)+2^w E_w(4z).
\]
This follows from the leading term of $p_s(m)$, after extracting $(nm)^{2p}$, followed by the parity split of Section~2. Each $G_s$ is modular, its constant term is $-(-1)^s$, and its first coefficient is the nonzero first coefficient of $E_w$. Their weights $w$ are distinct. Lemma~\ref{deriv} with $r=2p$, followed by top-weight reduction and specialization, proves the result.
\end{proof}

\clearpage
\section{Uniform algebraic independence at every positive degree}\label{global}
\begin{theorem}\label{full}
Let $p\ge1$, $1\le s_1<s_2<s_3$, and $W_1,W_2,W_3\in K[X]$ have degree exactly $p$. Then the three numbers
\[
 \sum_{n\ge1}\frac{W_i(n^2)}{F_n^{2s_i}},\qquad 1\le i\le3,
\]
are algebraically independent. The assertion also holds upon replacing $F_n$ by $U_n=(\alpha^n-\beta^n)/(\alpha-\beta)$ for any algebraic $\alpha\beta=-1$ with $0<|\beta|<1$.
\end{theorem}
The proof treats all possible placements of the denominator indices relative to $p$. We work with the highest numerator and denominator monomials. Lower terms have strictly smaller weight as in Section~\ref{quartic}. Put
\[
 d=\min(2p,2s-1),\quad 2k=|2p-2s+1|+1,\quad
 A=2p,\quad B_s=2s-1 .
\]
Call $s\le p$ the lower range and $s\ge p+1$ the upper range. In the lower range $s=p+1-k$, $d=2p+1-2k$; in the upper range $s=p+k$, $d=2p$. In both ranges $d(2k+d-1)=A B_s$ and $2k+2d=A+B_s+1$.

Write $T_k=E_{2k}(z+1/2)$ and $J_k=E_{2k}(2z)$. The top component is a nonzero multiple of $D^d g_s$ with
\begin{equation}\label{symboltable}
 g_s=\begin{cases}
 T_k & s\text{ even},\\
 -T_k+4^kJ_k & s\text{ odd, lower range},\\
 -T_k+2J_k & s\text{ odd, upper range}.
 \end{cases}
\end{equation}
For even $s$, the sign in \eqref{raw} is $(-1)^{nm}$. For odd $s$ it is $-(-1)^{nm}+2\boldsymbol1_{2\mid n}$. After $D^d$ is extracted, the residual power $2k-1$ falls on $n$ in the lower range or on $m$ in the upper range. Restricting $n$ to even indices therefore contributes $4^k J_k$ or $2J_k$, respectively. This proves \eqref{symboltable}. When $k=1$, $Lg_s$ is constant; otherwise $Lg_s=0$. In weight two the anomaly coefficient $Lg_s/12$ equals $1$ for $T_1$, $1$ for $-T_1+4J_1$, and $0$ for $-T_1+2J_1$.

For $s=1$ the middle entry in the differential row for $Dg_1$ has nonzero constant $2p(4^p-1)$, and the other two entries have zero constant. For $s>1$ all three entries start at $q$; after normalization their first coefficient is
\[
 v_s=(1,A B_s,A+B_s+1).
\]
Here we used $[q]g_s\ne0$, $d\ge2$ for $s>1$, and the commutator formula in Section~\ref{quartic}. Any two distinct slopes in the last column therefore give a nonzero $q^2$ determinant when $s=1$ is present. We may assume $s_i\ge2$.

\clearpage
\subsection*{Mixed parity}
Work on the normal subgroup $\Gamma(4)$ and put $\gamma=\left(\begin{smallmatrix}1&0\\2&1\end{smallmatrix}\right)$. The slash action by $\gamma$ on homogeneous modular forms extends to a graded ring automorphism of the quasimodular ring by fixing $E_2$. It commutes with $D,L,W$. For $D$, this follows on modular generators by covariance of the Serre derivative $Df-(w/12)E_2 f$ and on $E_2$ from $DE_2=(E_2^2-E_4)/12$; the other two commutations follow on generators. This is equivalently slash on almost holomorphic completions followed by taking the holomorphic part \cite{NSZ}.

Under this automorphism $E_{2k}(z),E_{2k}(2z)$ are fixed, and $E_{2k}(4z)$ goes to $4^{-k}E_{2k}(z+1/2)$. Indeed $4\gamma z=2-1/(z+1/2)$, and $2\gamma z=(2z)/(2z+1)$. The same calculation with $E_2^*=E_2-3/(\pi\operatorname{Im}z)$ gives the weight-two formula upon taking holomorphic parts. The half-translation identity of Section~2 shows that $T_k$ goes to $4^k E_{2k}(4z)$. Consequently the first positive Fourier exponent of the transformed $g_s$ is $\lambda_s=4$ if $s$ is even and $\lambda_s=2$ if $s$ is odd. The first coefficient is nonzero in every case.

For $s\ge2$, entries of the transformed differential row start at $q^{\lambda_s}$. After scaling, this first row coefficient is
\[
 v_\lambda(s)=\left(\lambda,\frac{A B_s}{\lambda},A+B_s+1\right).
\]
If the triple has mixed parity, two rows have the same $\lambda$ and span the plane with normal $N_\lambda=(A(A+1)/\lambda,\lambda,-A)$. They are independent since their $s$ are distinct. For the remaining row, with $\mu\ne\lambda$,
\[
 N_\lambda\cdot v_\mu(s)=A(\mu/\lambda-1)\bigl(A+1-B_s\lambda/\mu\bigr).
\]
This is nonzero: $A+1$ and $B_s$ are odd integers, whereas $\lambda/\mu$ is $2$ or $1/2$. The leading coefficient of the differential determinant proves independence.

\subsection*{A test using the coefficient of highest degree in $E_2$}
For a modular $g$ of positive weight $2k$, the leading term of $D^d g$ as a polynomial in $E_2$ is a nonzero constant times $g E_2^d$, by the Serre derivative. For three such modular bases $g_i$, let $e$ be a nonzero vector spanning the integer nullspace of the rows $(k_i)_i,(d_i)_i$ when these rows have rank two. If $\prod_i g_i^{e_i}$ is nonconstant then the three $D^{d_i}g_i$ are algebraically independent. Indeed, in a relation on the leading terms, group monomials by $E_2$ degree and modular weight. In a common group their exponent differences are multiples of primitive $e$, which would give a polynomial equation with constant coefficients for the nonconstant meromorphic quotient. This is Wang's derivative criterion \cite[Proposition 5.5]{Wang2025}. Independence of leading terms implies independence of the original polynomials by taking maximal weighted degree in a putative relation. Any two nonzero modular forms of weights $2k,2l$ are independent whenever their weight-zero quotient $g^l/h^k$ is nonconstant, by the same grading argument.

\clearpage
\subsection*{Three odd denominator indices}
First suppose the three $g_i$ are modular. In a common range the $k_i$ are distinct and $(k_i),(d_i)$ have rank two. Across ranges, two indices share one range, so the rank is again two. In the lower range the transformed odd $g_s$ is $4^k(E_{2k}(2z)-E_{2k}(4z))$. It has vanishing order two. After dividing by its first term, its expansion is $1+2^{2k-1}q^2+O(q^4)$, by the Fourier coefficient $1+2^{2k-1}$ of the normalized divisor sum at two. The upper odd $g_s$ has constant $1$ at infinity and transformed constant $2-4^k\ne0$.

If all three indices lie in the lower range then $\sum e_i=\sum e_i k_i=0$ since $d_i=2p+1-2k_i$. After transformation, the first surviving test for constancy of the quotient is $\sum e_i2^{2k_i-1}$. This is nonzero by strict convexity in $k$, since the $k_i$ are distinct. If all three are upper, the same two sums vanish since $d_i=2p$. Comparing absolute constants at the two cusps gives $\sum e_i\log(4^{k_i}-2)$, which is nonzero by strict concavity of $x\mapsto\log(4^x-2)$ on $x\ge1$.

With one lower and two upper indices the exponent on the lower form is nonzero (the two upper $k$ differ and both have $d=2p$). Its zero at the transformed cusp proves nonconstancy. With two lower and one upper it suffices that the sum of the two lower exponents is nonzero. If that sum vanished, inserting $d_i=2p+1-2k_i$ on the lower pair into $\sum e_i d_i=\sum e_i k_i=0$ would give $(2p+2k_{\rm upper})e_{\rm upper}=0$, then both lower exponents would also vanish. This finishes the modular cases.

There can be one anomalous odd base, namely $s=p$ with $p$ odd. Its highest $E_2$ term after differentiation is a nonzero constant times $E_2^{d+1}$, from its nonzero $E_2$ coefficient and Ramanujan's derivative formula. It suffices that the other two modular bases be independent: the three leading monomials, in $E_2$ and those two bases, then have linearly independent exponent vectors. For two lower bases the orders two at the transformed cusp and their unequal weights give a nonconstant weight-zero quotient. For a lower and an upper base, their orders two and zero suffice. For two upper bases compare normalized cusp constants: $(4^k-2)^{1/k}$ strictly increases with $k\ge1$, since its logarithm is $\log4+k^{-1}\log(1-2/4^k)$. The negative correction increases strictly. This proves all odd triples.

\subsection*{Three even denominator indices}
Translation by $1/2$ reduces the bases to $E_{2k}(z)$ and commutes with $D$. The $k_i$ are pairwise distinct, since equal $k$ in opposite ranges would give $s=p+k$ and $s'=p+1-k$, whose parities differ. For an even $s$ the range is determined by $k$: if $k\equiv p\pmod2$ then $(s,d)=(p+k,2p)$, otherwise $(s,d)=(p+1-k,2p+1-2k)$. These formulas always refer to indices $s\ge2$ present in the triple.

\clearpage
If $k_i\ge2$, the preceding modular criterion applies with rank two. A constant monomial quotient must, after cancellation, involve only $k\in\{2,3,4,5,7\}$, by Griffin and coauthors \cite[Theorem 1.1]{Griffin}. A quotient supported on just two distinct weights would require the pair $k=2,4$, since $E_8=E_4^2$ is the only power relation between distinct Eisenstein series in this list. Such support would require $d_2/2=d_4/4$. The indices $2,4$ belong to the same range, where $d/k$ is strictly monotone, so this is impossible. Hence if a constant quotient is possible, all three $k$ lie in that five-element set.

Write $E_{2k}=E_4^{a_k}E_6^{b_k}$ with exponent pairs
\[
 (a_k,b_k)=(1,0),(0,1),(2,0),(1,1),(2,1)
 \quad\text{for }k=2,3,4,5,7.
\]
The following determinants of rows $(a_k,b_k,d_k)$ use the range assignment just given. A nonzero entry gives independent leading monomials $E_4^{a_k}E_6^{b_k}E_2^{d_k}$.
\[
\begin{array}{c|rr@{\qquad}c|rr}
 k\text{ triple}&p\text{ even}&p\text{ odd}&k\text{ triple}&p\text{ even}&p\text{ odd}\\\hline
 (2,3,4)&-2p&-2p-1 &(2,5,7)&-2(p+2)&3-2p\\
 (2,3,5)&-2(p+2)&3-2p &(3,4,5)&2(p+4)&2p-7\\
 (2,3,7)&-4(p+2)&-2(2p-3)&(3,4,7)&4(p+4)&2(2p-7)\\
 (2,4,5)&2p&2p+1 &(3,5,7)&0&0\\
 (2,4,7)&2p&2p+1 &(4,5,7)&-2(p+4)&7-2p
\end{array}
\]
If $k=1$ occurs, the corresponding translated base is $E_2$ itself. Its derivative has leading term a nonzero scalar times $E_2^{d+1}$. The other two bases are independent unless their $k$ values are $2,4$, again by \cite{Griffin}. Only $(1,2,4)$ and $(3,5,7)$ now require a separate coefficient.

Use the original half-translated functions at infinity, before translating the bases to $E_{2k}(z)$. Their ratio of Fourier coefficients $[q^2]/[q]$ is $\rho_k=-1-4^k/2$. As in \eqref{cert}, the first two Fourier vectors of the differential row give a nonzero $q^4$ determinant provided the rows $(1,s,h_s)$ have nonzero determinant, with
\[
 h_s=2^{d-2p}(4p+3-2s)(-1-4^k/2).
\]
Indeed $N=(A(A+1),1,-A)$ annihilates $v_s$. Its scalar product with the next Fourier vector, divided by $[q]g_s$, is $2^{d-1}A(2A+3-2s)\rho_k$. This is a fixed nonzero multiple of $h_s$ for fixed $p$. Substitution of $(s,d)$ gives
\[
\begin{array}{c|cc}
 k\text{ triple}&p\text{ even}&p\text{ odd}\\\hline
 (1,2,4)&-15(30p-83)&15(30p+7)/64\\
 (3,5,7)&15(30p+97)/4096&-960(30p-173).
\end{array}
\]
None vanishes at a positive integer $p$. This completes the even case and the proof of functional independence in Theorem~\ref{full}. Wang's specialization at $q_0$ gives the result for Fibonacci numbers. For the stated $U_n$, Binet's expansion uses the same double series at $q=\beta^2$ and replaces $5^s$ by $(\alpha-\beta)^{2s}$. This algebraic factor is nonzero, and $0<|q|<1$, as required.

\clearpage
\paragraph{Scope.}
For the four ordinary and alternating Fibonacci and Lucas families, Theorem~\ref{rank} classifies all algebraic-coefficient linear relations on polynomial weight spaces. For ordinary Fibonacci sums of any common positive numerator degree, Theorem~\ref{full} classifies triples as algebraically independent. The full level-four field still has transcendence degree three. The first-power reciprocal Fibonacci constant $\sum F_n^{-1}$ has an odd denominator exponent and lies outside these statements. Stein's original unweighted conjecture has a published proof; the present polynomial-weight extension and common-degree theorem are research results in this manuscript, whose publication priority remains unverified.

\begin{thebibliography}{9}
\bibitem{Griffin} T. Griffin, N. Kenshur, A. Price, B. Vandenberg-Daves, H. Xue, and D. Zhu, Monomials of Eisenstein series, \emph{J. Number Theory} \textbf{219} (2021), 445--459. \href{https://doi.org/10.1016/j.jnt.2020.09.002}{DOI}.
\bibitem{ET} C. Elsner and N. Technau, On linear relations for Dirichlet series formed by recursive sequences of second order, \emph{J. Aust. Math. Soc.} \textbf{110} (2021), 406--430. \href{https://doi.org/10.1017/S1446788719000491}{DOI}. \href{https://arxiv.org/abs/1805.03003}{arXiv:1805.03003}.
\bibitem{Wang2025} W. Wang, Evaluation of reciprocal sums of hyperbolic functions using quasimodular forms, \emph{J. Math. Anal. Appl.} \textbf{549} (2025), 129526. \href{https://doi.org/10.1016/j.jmaa.2025.129526}{DOI}. Public text used here: \href{https://arxiv.org/html/2312.16929v2}{arXiv:2312.16929v2}, especially Theorem 5.3, Lemma 5.4, and Proposition 5.10.
\bibitem{Wang2024} W. Wang, Algebraic independence of values of quasi-modular forms, \emph{J. Number Theory} \textbf{255} (2024), 85--97. \href{https://doi.org/10.1016/j.jnt.2023.08.010}{DOI}.
\bibitem{NSZ} K. Nagatomo, Y. Sakai, and D. Zagier, Modular linear differential operators and generalized Rankin--Cohen brackets. \href{https://people.mpim-bonn.mpg.de/zagier/files/preprints/MLDO.pdf}{Authors' manuscript}, Section 8, equation (40).
\end{thebibliography}
\end{document}
